\section*{الفصل \ref{ch:b1:taylor} --- صيغ تايلور والنشور
المقاربة}

\begin{solution}{exo:b1:taylor:1}
\[
\eu^{2x} = 1 + 2x + 2x^2 + \frac{4x^3}{3} + o(x^3);
\qquad
\ln(1 - x) = -x - \frac{x^2}{2} - \frac{x^3}{3} - \frac{x^4}{4} +
o(x^4);
\]
\[
\sqrt{1+x} = 1 + \frac x2 - \frac{x^2}{8} + \frac{x^3}{16} + o(x^3);
\qquad
\frac{1}{1 + x^2} = 1 - x^2 + x^4 - x^6 + o(x^6),
\]
والأخير بتعويض $-x^2$ في النشر الهندسي.
\end{solution}

\begin{solution}{exo:b1:taylor:2}
$\eu^x - 1 - x = \frac{x^2}{2} + o(x^2)$: فالنهاية $\dfrac12$.

$\cos x = 1 - \frac{x^2}{2} + \frac{x^4}{24} + o(x^4)$ و
$\sqrt{1 - x^2} = 1 - \frac{x^2}{2} - \frac{x^4}{8} + o(x^4)$:
والفرق $\frac{x^4}{24} + \frac{x^4}{8} = \frac{x^4}{6} + o(x^4)$:
فالنهاية $\dfrac16$.

$\ln(1+x) - \sin x = \bigl(x - \frac{x^2}{2}\bigr) - x + o(x^2) =
-\frac{x^2}{2} + o(x^2)$: فالنهاية $-\dfrac12$.
\end{solution}

\begin{solution}{exo:b1:taylor:3}
$\frac{1}{1+t^2} = 1 - t^2 + t^4 + o(t^5)$؛ وبالمكاملة من $0$ إلى
$x$ (\cref{met:b1:taylor:compute}~(5)):
\[
\arctan x = x - \frac{x^3}{3} + \frac{x^5}{5} + o(x^5)\ \ (\text{زوجية
}o(x^6)\text{، بالفردية}).
\]
$(1 - t^2)^{-1/2} = 1 + \frac{t^2}{2} + \frac38 t^4 + o(t^5)$ (بنشر ثنائي
الحدّ مع $\alpha = -\frac12$ و $x = -t^2$:
$\binom{-1/2}{2} = \frac{(-\frac12)(-\frac32)}{2} = \frac38$)؛
وبالمكاملة:
\[
\arcsin x = x + \frac{x^3}{6} + \frac{3x^5}{40} + o(x^5) .
\]
\end{solution}

\begin{solution}{exo:b1:taylor:4}
تايلور--لاغرانج (\cref{thm:b1:taylor:lagrange}) من أجل $\exp$ عند $a =
0$ مع $x = 1$: \omterm{def:b1:derivative:def}{فالمشتقة} $(n+1)$ هي $\eu^t \leq \eu < 3$ على
$\intcc{0}{1}$، ومنه
\[
\Bigl|\eu - \sum_{k=0}^{n} \frac{1}{k!}\Bigr| \leq
\frac{3}{(n+1)!} .
\]
ومن أجل $6$ أرقام عشرية مضبوطة، نريد $\frac{3}{(n+1)!} < 5\times 10^{-7}$،
أي $(n+1)! > 6\times 10^{6}$: ولأن $10! = 3\,628\,800$ و $11! =
39\,916\,800$، يكفي $n + 1 = 11$، أي $n = 10$.
\end{solution}

\begin{solution}{exo:b1:taylor:5}
$n\ln\bigl(1 + \frac1n\bigr) = n\Bigl(\frac1n - \frac{1}{2n^2} +
\frac{1}{3n^3} + o\bigl(\frac{1}{n^3}\bigr)\Bigr) = 1 - \frac{1}{2n}
+ \frac{1}{3n^2} + o\bigl(\frac{1}{n^2}\bigr)$. وبأخذ الأسّي، مع
$u = -\frac{1}{2n} + \frac{1}{3n^2}$ و $\eu^u = 1 + u + \frac{u^2}2
+ o(u^2)$:
\[
\Bigl(1 + \frac1n\Bigr)^n = \eu\cdot \eu^{u}
= \eu\Bigl(1 - \frac{1}{2n} + \frac{1}{3n^2} + \frac{1}{8n^2}
+ o\Bigl(\frac{1}{n^2}\Bigr)\Bigr)
= \eu\Bigl(1 - \frac{1}{2n} + \frac{11}{24n^2} +
o\Bigl(\frac{1}{n^2}\Bigr)\Bigr).
\]
فالنهاية $\eu$؛ والخطأ $\sim \dfrac{\eu}{2n}$: أي بطيء (رقم واحد لكل
زيادة عشرية في $n$).
\end{solution}

\begin{solution}{exo:b1:taylor:6}
$f(x) = x^2 - x^4 = x^2(1 + o(1))$: فأول حدّ $x^2$، و $p = 2$ زوجي،
والمعامل $> 0$: أي قيمة صغرى موضعية عند $0$ (وليست شاملة: $f(2) = -12$).

$g(x) = x^3 + x^5$: فأول حدّ $x^3$، و $p$ فردي: فلا قيمة حدّية؛ و $g$
تعبر مماسها (الأفقي): أي انعطاف عند $0$.

$h = \exp$ عند $a = 1$: $h(x) = \eu + \eu(x-1) + \frac{\eu}{2}(x-1)^2
+ o((x-1)^2)$؛ والفرق مع المماس هو
$\frac{\eu}{2}(x-1)^2 + o(\cdot) > 0$ موضعيًا. \omterm{def:b1:logic:inj}{وشاملًا}: $\eu^x -
\eu x \geq 0$ من أجل كل $x$ بالتحدّب
(\cref{thm:b1:derivative:convexchar}~(3)): فيقع المنحنى فوق كل
مماس، والمساواة عند نقطة التماسّ فقط.
\end{solution}

\begin{solution}{exo:b1:taylor:7}
من أجل $x \to +\infty$:
\[
f(x) = x\Bigl(1 + \frac1x\Bigr)^{1/3}
= x\Bigl(1 + \frac{1}{3x} - \frac{1}{9x^2} +
o\Bigl(\frac{1}{x^2}\Bigr)\Bigr)
= x + \frac13 - \frac{1}{9x} + o\Bigl(\frac1x\Bigr):
\]
فالمقارب $y = x + \frac13$، والمنحنى تحته بجوار $+\infty$. وعندما $x
\to -\infty$، يصحّ الحساب نفسه (لأن الجذر التكعيبي
معرَّف من أجل كل الأعداد الحقيقية، و $\frac1x \to 0$): فالمقارب نفسه $y = x +
\frac13$، لكن الآن $-\frac{1}{9x} > 0$: فالمنحنى \emph{فوق} المستقيم.
\end{solution}

\begin{solution}{exo:b1:taylor:8}
$u_n = \sqrt n\bigl(\sqrt{1 + \tfrac1n} - 1\bigr) = \sqrt
n\bigl(\frac{1}{2n} + o(\frac1n)\bigr) \sim \dfrac{1}{2\sqrt n}$.

$v_n = \ln\bigl(1 + \frac1n\bigr) \sim \dfrac 1n$.

$w_n$: مع $h = \frac1n \to 0$، $\sin h - \tan h = \bigl(h -
\frac{h^3}{6}\bigr) - \bigl(h + \frac{h^3}{3}\bigr) + o(h^3) =
-\frac{h^3}{2} + o(h^3)$، ومنه $w_n \sim -\dfrac{1}{2n^3}$.
\end{solution}

\begin{solution}{exo:b1:taylor:9}
تايلور--لاغرانج من الرتبة $1$ حول $x$، على الجهتين:
\[
f(x + h) = f(x) + h f'(x) + R_+,\quad
f(x - h) = f(x) - h f'(x) + R_-,
\qquad \abs{R_\pm} \leq \frac{h^2}{2} \sup \abs{f''} .
\]
وبالطرح: $f(x+h) - f(x-h) = 2h f'(x) + (R_+ - R_-)$، ومنه
\[
\abs{f'(x)} \leq \frac{\abs{f(x+h) - f(x-h)}}{2h} +
\frac{h}{2}\sup_{\intcc{x-h}{x+h}}\abs{f''} .
\]
ومع الحواصر الشاملة: $\abs{f'(x)} \leq \frac{M_0}{h} + \frac{M_2
h}{2}$ من أجل كل $h > 0$. ويُصغَّر الطرف الأيمن عند $h =
\sqrt{2M_0/M_2}$ (حيث تنعدم \omterm{def:b1:derivative:def}{المشتقة})، بالقيمة $\sqrt{2M_0M_2}$
--- ومنه $\abs{f'} \leq \sqrt{2M_0M_2}$. (وإذا كان $M_2 = 0$، فدع $h \to
\infty$: فيكون $f' = 0$، وهذا متسق.)
\end{solution}

\begin{solution}{exo:b1:taylor:10}
على $\intoo{0}{\pi}$: $0 < \sin u < u$، ومنه فالمتتالية $(u_n)$ متناقصة
قطعًا وموجبة، ومنه متقاربة؛ والنهاية نقطة صامدة
للمقدار $\sin$ في $\intcc{0}{\pi}$، ويفرض $\sin \ell = \ell$ أن $\ell =
0$ (لأن $\sin x < x$ من أجل $x > 0$).
\begin{enumerate}
  \item باستعمال $\sin u = u - \frac{u^3}{6} + o(u^3)$ عندما $u \to 0$:
        \[
        v_{n+1} - v_n = \frac{1}{\sin^2 u_n} - \frac{1}{u_n^2}
        = \frac{1}{u_n^2}\Bigl(\Bigl(1 - \frac{u_n^2}{6} +
        o(u_n^2)\Bigr)^{\!-2} - 1\Bigr)
        = \frac{1}{u_n^2}\Bigl(\frac{u_n^2}{3} + o(u_n^2)\Bigr)
        \longrightarrow \frac13 .
        \]
  \item وحسب \cref{exo:b1:seq:10}~(3) (تشيزارو من أجل الفروق)،
        $\frac{v_n}{n} \to \frac13$، أي $v_n \sim \frac n3$، أي
        $u_n^2 \sim \frac 3n$: ولأن $u_n > 0$،
        \[
        u_n \sim \sqrt{\frac{3}{n}} .
        \]
\end{enumerate}
\end{solution}

\begin{solution}{exo:b1:taylor:11}
بالمقامات المشتركة. النهاية الأولى:
\[
\frac{1}{x^2} - \frac{1}{\sin^2 x}
= \frac{\sin^2 x - x^2}{x^2\sin^2 x},
\qquad
\sin^2 x = \Bigl(x - \frac{x^3}{6} + o(x^4)\Bigr)^{\!2}
= x^2 - \frac{x^4}{3} + o(x^5) ,
\]
ومنه فالبسط هو $-\frac{x^4}{3} + o(x^4)$ بينما
المقام $\sim x^4$: فالنهاية هي $-\dfrac13$.

والثانية: $\dfrac1x - \dfrac{1}{\ln(1+x)} = \dfrac{\ln(1+x) -
x}{x\ln(1+x)}$؛ والبسط هو $-\frac{x^2}{2} + o(x^2)$، و
المقام $x\bigl(x + o(x)\bigr) \sim x^2$: فالنهاية هي
$-\dfrac12$.
\end{solution}

\begin{solution}{exo:b1:taylor:12}
على $I_k = \intoo{k\pi - \frac\pi2}{k\pi + \frac\pi2}$، تكون
للدالة $g(x) = \tan x - x$ \omterm{def:b1:derivative:def}{مشتقة} $\tan^2 x \geq 0$،
لا تنعدم إلا عند النقطة الوحيدة $k\pi$: ومنه فالدالة $g$ متزايدة
قطعًا على $I_k$ (\cref{cor:b1:derivative:monotone}~(2))،
بنهايتين $-\infty$ و $+\infty$ عند الطرفين: أي جذر واحد
بالضبط $x_k$. ومن أجل $k \geq 1$، $g(k\pi) = -k\pi < 0$، ومنه $x_k \in
\intoo{k\pi}{k\pi + \frac\pi2}$: فاكتب $x_k = k\pi + \frac\pi2 -
\varepsilon_k$ مع $\varepsilon_k \in \intoo{0}{\frac\pi2}$.
عندئذ
\[
x_k = \tan x_k = \tan\Bigl(\frac\pi2 - \varepsilon_k\Bigr)
= \frac{1}{\tan\varepsilon_k}
\quad\Longrightarrow\quad
\varepsilon_k = \arctan\frac{1}{x_k} ,
\]
باستعمال $\tan\varepsilon_k = \frac{1}{x_k}$ و $\varepsilon_k \in
\intoo{0}{\frac\pi2}$. ولأن $x_k \geq k\pi \to \infty$:
$\varepsilon_k \to 0$، و
\[
\varepsilon_k = \arctan\frac{1}{x_k} = \frac{1}{x_k} +
O\Bigl(\frac{1}{x_k^3}\Bigr)
= \frac{1}{k\pi + O(1)} + O\Bigl(\frac{1}{k^3}\Bigr)
= \frac{1}{k\pi} + O\Bigl(\frac{1}{k^2}\Bigr) ,
\]
ومنه $x_k = k\pi + \frac\pi2 - \frac{1}{k\pi} +
o\bigl(\frac1k\bigr)$.
\end{solution}

\begin{solution}{pb:b1:taylor:1}
\textbf{1.} $S_{2n+1} - S_{2n-1} = a_{2n} - a_{2n+1} \geq 0$ و
$S_{2n+2} - S_{2n} = a_{2n+2} - a_{2n+1} \leq 0$، بينما $S_{2n}
- S_{2n+1} = a_{2n+1} \to 0$: ومنه فالمتتاليتان $(S_{2n+1})$ و
$(S_{2n})$ متجاورتان، وتتقاربان إلى $S$ مشترك
(\cref{thm:b1:seq:adjacent}) مع $S_{2n+1} \leq S \leq S_{2n}$.
ومن أجل $n$ زوجي: يعطي $S_{n+1} \leq S \leq S_n$ أن $-a_{n+1} \leq S -
S_n \leq 0$؛ ومن أجل $n$ فردي: $0 \leq S - S_n \leq a_{n+1}$. وفي
الحالتين $\abs{S - S_n} \leq a_{n+1}$ ويحمل $S - S_n$ إشارة
$(-1)^{n+1}$، أي أول حدّ محذوف. والتناقص القطعي يجعل
كل متراجحة معروضة قطعية، وعلى الخصوص $0 < \abs{S -
S_n} < a_{n+1}$.

\textbf{2.} المقدار $a_k = \frac{1}{k!}$ يتناقص قطعًا إلى $0$:
ومنه ينطبق السؤال 1. وتايلور--لاغرانج
(\cref{thm:b1:taylor:lagrange}) من أجل $\exp$ بين $-1$ و $0$:
$\abs{\eu^{-1} - T_n} \leq \frac{1}{(n+1)!}$ (لأن \omterm{def:b1:derivative:def}{المشتقة}
$\eu^t$ هي $\leq 1$ هناك)، ومنه $T_n \to \eu^{-1}$، وتكون النهاية
$S$ في السؤال 1 \emph{هي} $\eu^{-1}$، بالحاصرين القطعيين
$0 < \abs{\eu^{-1} - T_n} < \frac{1}{(n+1)!}$.

\textbf{3.} بمكاملة المتطابقة على $\intcc{0}{1}$: يكون
الطرف الأيسر $\arctan 1 = \frac\pi4$ (بالمبرهنة الأساسية)، ويعطي
الحدّ $k$ المقدارَ $\frac{(-1)^k}{2k+1}$، و
\[
\rho_n = (-1)^{n+1}\int_0^1 \frac{t^{2n+2}}{1+t^2}\dd t,
\qquad
\abs{\rho_n} \leq \int_0^1 t^{2n+2}\dd t = \frac{1}{2n+3} .
\]

\textbf{4.} الخطأ على $\pi$ هو $4\abs{\rho_n} \leq
\frac{4}{2n+3}$: ويقتضي نزوله دون $10^{-6}$ أن $2n + 3 > 4\cdot10^6$،
أي نحو مليونَي حدّ. وفي الأثناء $4S_4 = 4\bigl(1 - \frac13 +
\frac15 - \frac17 + \frac19\bigr) = 4 \times 0.834921 =
3.339683$، أي على بعد $0.2$ تقريبًا من $\pi$: خمسة حدود، ولا حتى
رقم واحد.

\textbf{5.} بالمكاملة على $\intcc{0}{1}$: $\ln 2 =
\sum_{k=1}^{n} \frac{(-1)^{k-1}}{k} + (-1)^n R_n$ مع $R_n =
\int_0^1 \frac{t^n}{1+t}\dd t$؛ ومن $\frac12 \leq \frac{1}{1+t}
\leq 1$: $\frac{1}{2(n+1)} \leq R_n \leq \frac{1}{n+1}$. و
الخطأ محبوس بين مضاعفين للمقدار $\frac1n$: فستة
أرقام عشرية تكلّف نحو مليون حدّ.

\textbf{6.} بالمكاملة من $0$ إلى $x$: $\frac12\ln\frac{1 +
x}{1 - x} = \sum_{k=0}^{n} \frac{x^{2k+1}}{2k+1} + \int_0^x
\frac{t^{2n+2}}{1-t^2}\dd t$. وعند $x = \frac13$: $\frac{1 +
1/3}{1 - 1/3} = 2$، وعلى $\intcc{0}{\frac13}$، $\frac{1}{1 -
t^2} \leq \frac98$:
\[
\ln 2 = 2\sum_{k=0}^{n} \frac{(1/3)^{2k+1}}{2k+1} +
\tilde\rho_n,
\qquad
0 < \tilde\rho_n \leq \frac{9}{4}\cdot
\frac{(1/3)^{2n+3}}{2n+3} :
\]
فكل حدّ إضافي \omterm{def:b1:arith:divides}{يقسم} الخطأ على نحو $9$.

\textbf{7.} بالاستقراء: من أجل $n = 1$: $1 - \frac12 = \frac12 = H_2
- H_1$. والخطوة:
\[
H_{2n+2} - H_{n+1} = (H_{2n} - H_n) + \frac{1}{2n+1} +
\frac{1}{2n+2} - \frac{1}{n+1}
= (H_{2n} - H_n) + \frac{1}{2n+1} - \frac{1}{2n+2},
\]
وهو بالضبط تزايد المجموع المتناوب. و
$H_{2n} - H_n = \sum_{k=1}^{n}\frac{1}{n+k}$ هو مجموع ريمان
للمقدار \cref{ex:b1:integration:riemannexample}، وهو يتقارب إلى $\ln
2$: ومنه فالمجاميع الجزئية الزوجية في الطريق 1 \emph{هي} \omterm{thm:b1:integration:riemann}{مجاميع ريمان}
في الطريق 3.

\textbf{8.} $\sum_{k=1}^{6}\frac{(-1)^{k-1}}{k} = 0.61667$،
بخطأ $0.0765$؛ ويعطي الطريق 2 عند $n = 5$ المقدارَ $0.6931471$ بخطأ
$\leq \frac94\cdot\frac{(1/3)^{13}}{13} = 1.1\cdot10^{-7}$. و
السبب: يقوّم الطريق 1 متسلسلة اللوغاريتم عند النقطة الحدّية
$x = 1$، حيث تتلاشى الحدود مثل $\frac1k$؛ ويقوّم الطريق 2
عند $x = \frac13$، وهو عميق في الداخل، حيث يحمل كل حدّ
عاملًا جديدًا $\frac19$.

\textbf{9.} عشرة أرقام عشرية: نريد $\tilde\rho_n \leq
5\cdot10^{-11}$. وعند $n = 10$: $\frac94 \cdot
\frac{(1/3)^{23}}{23} = \frac94\cdot\frac{1.06\cdot10^{-11}}{23}
\approx 1.0\cdot10^{-12} < 5\cdot10^{-11}$: فأحد عشر حدًّا
تكفي.

\textbf{10.} $(5+\iu)^2 = 24 + 10\iu$، ثم $(5+\iu)^4 = (24 +
10\iu)^2 = 476 + 480\iu$؛ و $2(1+\iu)(239+\iu) = 2(238 +
240\iu) = 476 + 480\iu$: فهما متساويان. والعمد: $\arg(5 + \iu) =
\arctan\frac15$، ومنه فللطرف الأيسر \omterm{def:b1:complex:expi}{عمدة}
$4\arctan\frac15 \approx 0.79 \in \intoo{0}{\pi}$؛ وللطرف
الأيمن \omterm{def:b1:complex:expi}{عمدة} $\frac\pi4 + \arctan\frac{1}{239} \in
\intoo{0}{\pi}$. وعددان عقديان متساويان بعمدتين في
\omterm{prop:b1:reals:intervals}{الفترة} نفسها التي طولها $< 2\pi$:
\[
4\arctan\frac15 = \frac\pi4 + \arctan\frac{1}{239} ,
\]
وهي صيغة ماشان.

\textbf{11.} كامل $\frac{1}{1+t^2} = \sum_{k=0}^n (-1)^k
t^{2k} + \frac{(-1)^{n+1}t^{2n+2}}{1+t^2}$ من $0$ إلى $x$:
\[
\arctan x = \sum_{k=0}^{n}\frac{(-1)^k x^{2k+1}}{2k+1} +
r_n(x),
\qquad
\abs{r_n(x)} \leq \int_0^x t^{2n+2}\dd t =
\frac{x^{2n+3}}{2n+3} .
\]

\textbf{12.} الأخطاء: $16\,\frac{(1/5)^{11}}{11} = 3.0\cdot
10^{-8}$ و $4\,\frac{(1/239)^5}{5} \approx 1\cdot10^{-12}$:
والمجموع $< 5\cdot10^{-8}$. ويساوي المجموع المعروض
$3.14159268\dots$، ومنه $\pi = 3.1415926\dots$ مصادَقًا عليه حتى
$5\cdot10^{-8}$: أي سبعة أرقام عشرية من ستة حدود (خمسة عند
$\frac15$ واثنان عند $\frac1{239}$ بحساب سخيّ).

\textbf{13.} لايبنتز: خطأ $\sim \frac1n$، ومنه فكل رقم جديد
يضرب العمل في عشرة. ودالزل
(\cref{pb:b1:integration:1}، السؤال 22): خطأ $4^{1-5m}$،
أي نحو ثلاثة أرقام لكل خطوة، وكل خطوة \omterm{def:b1:poly:def}{كثير حدود} أثقل.
وماشان: نسبة خطأ $\frac{1}{25}$ لكل حدّ، أي نحو $1.4$ رقم
لكل حدّ، وكل حدّ قسمة واحدة. والعبرة: يتسلّم باقي
النشر من النوع الهندسي مثل $x^{2n}$، ومنه فتصغير $x$
يشتري الأرقام بكلفة \emph{ثابتة} لكل حدّ --- ومتطابقة
ماشان العقدية هي بالضبط آلة لتصغير $x$.

\textbf{14.} المقدار $a_k = \frac{1}{(2k)!}$ يتناقص قطعًا إلى $0$؛
وحسب السؤال 1 و \cref{prop:b1:taylor:standard} (بحاصر لاغرانج
كما في السؤال 2)، $\sum_{k \leq n}\frac{(-1)^k}{(2k)!}
\to \cos 1$ مع التأطير \emph{القطعي} $0 < \bigl|\cos 1 -
S'_n\bigr| < \frac{1}{(2n+2)!}$. افترض $\cos 1 = \frac pq$ و
خذ $2n \geq q$: عندئذ $(2n)!\,S'_n = \sum_{k\leq n} (-1)^k
\frac{(2n)!}{(2k)!} \in \Z$ و $(2n)!\,\frac pq \in \Z$، بينما
\[
0 < \Bigl|(2n)!\cos 1 - (2n)!S'_n\Bigr| <
\frac{(2n)!}{(2n+2)!} = \frac{1}{(2n+1)(2n+2)} < 1 :
\]
أي عدد صحيح غير معدوم قيمته المطلقة $< 1$. وهذا تناقض: ومنه $\cos
1 \notin \Q$.

\textbf{15.} بالمثل مع $a_k = \frac{1}{(2k+1)!}$،
وبالضرب في $(2n+1)!$ مع $2n + 1 \geq q$: $\sin 1 \notin
\Q$. ومع ذلك $\cos^2 1 + \sin^2 1 = 1 \in \Q$: فجداءات الأعداد
الصمّاء ومجاميعها قد تكون ناطقة --- فالصمم لا يمرّ عبر أيّ
عملية جبرية بالمجّان.

\textbf{16.} حاصر الذيل: من أجل $m > n$،
\[
\sum_{k=n+1}^{m} \frac{1}{(2k)!}
\leq \frac{1}{(2n+2)!}\Bigl(1 + \frac12 + \frac14 +
\dots\Bigr) \leq \frac{2}{(2n+2)!} ,
\]
لأن كل نسبة متتالية هي $\frac{1}{(2k+1)(2k+2)} \leq
\frac12$؛ والذيل موجب (لأن حدّه الأول كذلك). ومنه $0 <
\cosh 1 - \sum_{k\leq n}\frac{1}{(2k)!} < \frac{2}{(2n+2)!}$،
والضرب في $(2n)!$ مع $2n \geq q$ يحبس عددًا صحيحًا غير معدوم
في $\intoo{0}{1}$ مرة أخرى: ومنه $\cosh 1 \notin \Q$.

\textbf{17.} $\cos\frac1m = \sum_k
\frac{(-1)^k}{m^{2k}(2k)!}$: فالحدود تتناقص قطعًا إلى
الصفر، و $m^{2n}(2n)!\cdot\frac{1}{m^{2k}(2k)!} =
m^{2(n-k)}\frac{(2n)!}{(2k)!} \in \Z$ من أجل $k \leq n$. وإذا كان
$\cos\frac1m = \frac pq$، فاضرب التأطير القطعي في
$q\,m^{2n}(2n)!$: فيكون الخطأ محدودًا بالمقدار
$\frac{q}{m^2(2n+1)(2n+2)} < 1$ من أجل $n$ كبير: وهذا تناقض.
وأمّا من أجل $\frac ab$ مع $a \geq 2$: فإن إزالة المقامات تضرب
الذيلَ في $b^{2n}(2n)!$، لكن أول حدّ محذوف هو
$\frac{a^{2n+2}}{b^{2n+2}(2n+2)!}$، والجداء
$\frac{a^{2n+2}}{b^2(2n+1)(2n+2)}$ \emph{ينفجر}: فالبسط
$a^{2n+2}$ لم يعد يُزال، وينحشر الفخّ.
(والنتيجة صحيحة مع ذلك --- عبر آلة من نوع نيفن، لا هذه.)

\textbf{18.} عند $t = x$ ينعدم كل حدّ من $g$ إلا $f(x)
- f(x) = 0$: ومنه $g(x) = 0$؛ ويُختار $A$ بحيث $g(a) = 0$.
وبالاشتقاق، يتلسكب المجموع:
\[
g'(t) = -\frac{f^{(n+1)}(t)}{n!}(x - t)^n +
A\,\frac{(x-t)^n}{n!}
= \frac{(x-t)^n}{n!}\bigl(A - f^{(n+1)}(t)\bigr) .
\]
وتعطي رول على القطعة من $a$ إلى $x$ عددًا $c$ قطعًا بينهما
يحقق $g'(c) = 0$؛ ولأن $(x - c)^n \neq 0$: $A =
f^{(n+1)}(c)$. ونشرُ $g(a) = 0$ يعطي مساواة تايلور
بالباقي $\frac{f^{(n+1)}(c)}{(n+1)!}(x-a)^{n+1}$.

\textbf{19.} من أجل $x > 0$، يكون الباقي $\frac{\eu^{c}}
{(n+1)!}x^{n+1} > 0$: فالأسّي يتجاوز كلًّا من كثيرات حدود
تايلور له، قطعًا، عند كل رتبة. ومن أجل $x < 0$ تكون
إشارة الباقي إشارةَ $x^{n+1}$: فيكون $\eu^x$
\emph{فوق} \omterm{def:b1:poly:def}{كثير الحدود} من أجل $n$ فردي، و\emph{تحته} من أجل
$n$ زوجي --- أي جهتين متناوبتين، كما يبيّن أصلًا منحنيا $1 + x$ و $1
+ x + \frac{x^2}{2}$ إزاء $\eu^x$.

\textbf{20.} الخطأ الحقيقي: $\sin 0.5 - 0.4791667 =
2.59\cdot10^{-4}$، إزاء الحاصر $\frac{0.5^5}{120} =
2.60\cdot10^{-4}$: فيكاد يُبلغ (لأن الحدّ التالي يهيمن على
الذيل). ويونغ: من أجل النهايات والتحليل الموضعي، حيث لا حاجة إلى أيّ
ثابت. ولاغرانج: من أجل الأرقام العشرية المصادَق عليها. والمتناوب: حين
ينطبق، فالحاصر نفسه \emph{مع} اتجاه الخطأ
--- أي أفضل الثلاثة، لكنه أندرها.

\textbf{21.} \omterm{def:b1:continuity:continuous}{الاتصال} عند $0$: مع $u = \frac{1}{x^2} \to
+\infty$، $f(x) = \eu^{-u} \to 0 = f(0)$. \omterm{def:b1:derivative:def}{والمشتقة} عند $0$:
$\bigl|\frac{f(h)}{h}\bigr| = \sqrt u\,\eu^{-u} \to 0$
(\cref{prop:b1:functions:powerrules}): ومنه $f'(0) = 0$. ومن أجل $x \neq
0$، $f'(x) = \frac{2}{x^3}\eu^{-1/x^2}$: أي الشكل
$P_1\bigl(\frac1x\bigr)\eu^{-1/x^2}$ مع $P_1(X) = 2X^3$؛
وبالاستقراء، يعطي اشتقاق $P_k(\frac1x)\eu^{-1/x^2}$ المقدارَ
$P_{k+1}(X) = 2X^3 P_k(X) - X^2 P_k'(X)$، وهو \omterm{def:b1:poly:def}{كثير حدود}. وعندئذ
\[
\frac{f^{(k)}(h) - 0}{h} = v\,P_k(v)\,\eu^{-v^2}
\Big|_{v = 1/h} \longrightarrow 0
\]
(\omterm{def:b1:poly:def}{كثير حدود} إزاء $\eu^{-v^2}$، \omterm{prop:b1:functions:powerrules}{بمقارنة النمو} عند
$\pm\infty$): ومنه بالاستقراء $f^{(k)}(0) = 0$ من أجل كل $k$. فتنعدم كل
كثيرات حدود تايلور للمقدار $f$ عند $0$، ومع ذلك $f > 0$ خارج $0$:
فتايلور--يونغ مضبوط عند كل رتبة وأعمى بعد
الجرثومة. فالنشر معلومة موضعية فقط.

\textbf{22.} حسب السؤال 2، $0 < \abs{\eu^{-1} - T_n} <
\frac{1}{(n+1)!}$، قطعًا. فإذا كان $\eu^{-1} = \frac pq$، فخذ $n
\geq q$ واضرب في $n!$: فيكون $n!\,T_n \in \Z$ و $n!\frac pq
\in \Z$، ومنه يكون لعدد صحيح غير معدوم قيمة مطلقة $<
\frac{n!}{(n+1)!} = \frac{1}{n+1} < 1$: وهذا تناقض. ومنه
$\eu^{-1} \notin \Q$، و $\eu = \frac{1}{\eu^{-1}}$ أصمّ
كذلك. والبراهين الثلاثة: \omterm{thm:b1:seq:adjacent}{المتتاليتان المتجاورتان} بحصر
$q!\,\eu$ (\cref{exo:b1:seq:9})؛ والعلاقة التراجعية التكاملية $A_n =
\eu - nA_{n-1}$ (\cref{pb:b1:integration:1})؛ والتأطير
المتناوب (هنا). أي فخّ واحد وثلاث شهادات صغر.

\textbf{23.} اجمع المجاميع الجزئية أزواجًا: $\sum_{k=1}^{2n}
(-1)^k b_k = E_n - O_n$ مع $E_n = \sum_{j=1}^{n}
\frac{1}{4j^2}$، وهو محدود (بحاصر التلسكوب في
\cref{ex:b1:seq:basel}، $E_n \leq \frac12$)، و $O_n =
\sum_{j=1}^{n}\frac{1}{2j-1} \geq \frac12 H_n \to +\infty$
(\cref{exo:b1:seq:5}): ومنه تؤول المجاميع الجزئية إلى $-\infty$.
واستُعملت الرتابة في السؤال 1 بالضبط حيث احتاج $S_{2n+1} -
S_{2n-1} = a_{2n} - a_{2n+1}$ إلى إشارة: فبدون التناقص،
ليس على المتتاليتين الجزئيتين الزوجية والفردية أن تكونا رتيبتين، و
ينهار التجاور.

\textbf{24.} $(2+\iu)(3+\iu) = 5 + 5\iu = 5(1+\iu)$؛ وبأخذ
العمد (وكلها في $\intoo{0}{\frac\pi2}$): $\arctan\frac12 +
\arctan\frac13 = \frac\pi4$. وكلفة المتسلسلة من أجل ستة أرقام عشرية:
الخطأ $\leq 4\bigl(\frac{(1/2)^{2n+3}}{2n+3} +
\frac{(1/3)^{2n+3}}{2n+3}\bigr)$، وهو عند $n = 10$ يساوي
$\approx 2\cdot10^{-8} < 5\cdot10^{-7}$: أي أحد عشر حدًّا. والترتيب:
أفضل من لايبنتز بهامش أسّي، وخلف ماشان
(الذي نقطته المهيمنة $\frac15$ أصغر من $\frac12$):
أي نحو $0.6$ رقم لكل حدّ إزاء $1.4$ عند ماشان.

\textbf{25.} (أ) من أجل $a_k \to 0$ متناقصة، تتقارب المجاميع
الجزئية المتناوبة مع $\abs{S - S_n} \leq a_{n+1}$ ويحمل
الخطأ إشارة أول حدّ محذوف. (ب) والمتطابقة المنتهية
ذات الباقي الصريح يمكن \emph{تقويمها وحصرها
عند نقطة مختارة}، بينما لا تعد \omterm{def:b1:logic:statement}{عبارة} النهاية إلا بتقارب
في النهاية --- وتحتاج المصادقة إلى الأولى.
(ج) والمستخرج: $\pi$ حتى $5\cdot10^{-8}$ بماشان، و $\ln 2$ حتى
عشرة أرقام عشرية بمتسلسلة $\frac13$، وصمم $\cos 1$ و
$\sin 1$ و $\cosh 1$ و $\cos\frac1m$ و $\eu^{-1}$، و
مساواة تايلور--لاغرانج، وتحذير الدالة المنبسطة. (د)
ويرقّي \cref{ch:b1:series} السؤالَ 1 إلى محك المتسلسلات
المتناوبة، ويفصل التقارب بإطلاق عن التقارب الشرطي،
وتعرض مسألة نهاية الأسبوع فيه دراما إعادة الترتيب التي
تكون فيها المتسلسلة التوافقية المتناوبة في الطريق 1 الشاهدَ
النجم.
\end{solution}
